% FILE: 01_research_question.tex
% Project: schemas
% Thesis Role: data-schema-definitions-for-ma-and-receipt-surfaces

\section{Research Question}
\label{sec:schemas-research-question}

\subsection{Problem Statement}
\label{sec:schemas-research-question-problem}

Process intelligence claims made in M\&A due diligence contexts carry financial consequences.
A claim that annual EBITDA will increase by \$1,250,000 via purchase order rework reduction
(as asserted in \texttt{rec\_ebitda\_rework\_001.json}) is not an opinion---it is a quantified
projection that a board will rely upon in approving or rejecting a transaction. The problem is
that no structural contract previously existed between the process mining execution layer and the
board presentation layer. Claims could be assembled, formatted, and presented without any
mechanism to verify that the underlying event log hash, process model hash, conformance metric,
and cryptographic signature were present and structurally valid.

This gap is not merely an engineering inconvenience. Under the Van der Aalst doctrine applied in
this research program, a model-vs-log mismatch is a \emph{defect}, not a discrepancy. By
extension, a claim that cannot be structurally validated against a schema is a defect in the
evidence chain, regardless of whether the underlying numbers are correct.

The research question therefore is: \emph{what is the minimal, mechanically enforceable structural
contract that makes a process intelligence receipt admissible as board-level M\&A evidence?}

\subsection{Old-Regime Assumption Challenged}
\label{sec:schemas-research-question-old-regime}

The old-regime assumption is that process mining outputs---fitness scores, conformance verdicts,
variant analyses---are self-evidencing. Under this assumption, a practitioner produces a report
that asserts a fitness value of 0.982 and a precision of 0.945, and the board accepts those
numbers because the practitioner is credentialed and the tool is reputable. There is no structural
requirement that the assertion be linked to a specific event log by a cryptographic hash, that the
process model be identified by content address, or that the entire payload be signed by a verifiable
key.

This assumption fails in adversarial contexts. An auditor representing a counterparty in an
M\&A transaction cannot re-execute the analysis unless the event log is identified by hash and the
query is specified by URI. A court cannot evaluate the integrity of the claim without a signature.
A seller cannot defend a claim against a buyer's challenge without a receipt that is independent
of the seller's own assertions.

The schemas project challenges this assumption directly: it asserts that a receipt without a
structural contract is not a receipt---it is an informal note.

\subsection{Hypothesis}
\label{sec:schemas-research-question-hypothesis}

The hypothesis is that a JSON Schema with \texttt{patternProperties} keyed on UUIDv4 identifiers,
requiring BLAKE3 hashes of both the event log and the process model, a typed query definition with
an engine enum constraining execution to \texttt{wasm4pm} or \texttt{pm4py}, a verification status
enum, and a 128-hex Ed25519 signature, constitutes a sufficient structural contract for
board-admissible process intelligence receipts.

Sufficiency here means: any receipt conforming to this schema carries enough information for an
independent auditor to re-execute the query on the identified event log, compare the result to the
stored verification metrics, verify the cryptographic signature against a known public key, and
render a conformance verdict without relying on the original analyst's testimony.

This hypothesis is supported by the five sample receipts in the \texttt{receipts/} directory, each
of which conforms to the schema by inspection and carries sufficient content for independent
re-execution---subject to the caveat that no automated CI enforcement existed at time of assessment.

\subsection{Success Criteria}
\label{sec:schemas-research-question-success}

Success is defined across three levels:

\begin{enumerate}
  \item \textbf{Structural completeness}: Every required field is present in the schema, and the
    field patterns (64-hex BLAKE3, 128-hex Ed25519, UUIDv4) are tight enough to reject malformed
    values without accepting invalid content.
  \item \textbf{Receipt conformance}: All five production receipts (\texttt{rec\_ebitda\_rework\_001.json}
    through \texttt{rec\_residual\_standard\_005.json}) conform to the schema, demonstrating that
    the contract is practicable for real M\&A deck assertions.
  \item \textbf{Pipeline integration}: The schema is referenced by the M\&A deck rendering authority
    assessment (dated 2026-06-01) and by the slide-to-receipt map, establishing that the schema is
    the governing contract for the full seven-step rendering pipeline---not a standalone document.
\end{enumerate}

Levels 1 and 2 are met as of the PARTIAL status assessment. Level 3 is met at the documentation
level. The remaining gap is Level~0: no automated schema validation exists at commit time, meaning
the contract is not mechanically enforced in continuous integration.
